\section{Introduction}
\label{sec:intro}

The deep ocean is the largest habitat on Earth and the least catalogued. Below
the abyssal boundary at 4,000\,m, and in the hadal trenches that reach
11,000\,m, most expeditions still return with species new to science
\cite{ramirez2010deep}. The question that motivates this paper is a natural
one for the age of satellites and machine learning: \emph{could we find those
species from above --- from orbit, or at least from a distance --- and could an
AI recognise an animal that no one has ever described?}

The first half of the question is a physics question, and we answer it before
building anything. Seawater is opaque to almost everything a satellite can
emit or receive. Sunlight is gone within a few hundred metres, radar and radio
are absorbed within centimetres to metres, and gravity or magnetic anomalies
respond to rock, not to animals. What a satellite \emph{can} see is the
surface: temperature fronts, chlorophyll, and --- through the sea-surface
height that gravity imposes --- the bathymetry of seamounts and trenches where
life concentrates. Section~\ref{sec:budget} quantifies each modality with
closed forms: Beer--Lambert attenuation for light, the skin depth for
electromagnetic waves, the passive sonar equation with Thorp absorption and
Wenz ambient noise for sound, and a shed--decay--transport plume for
environmental DNA. One modality survives the trip: low-frequency sound. The
eyes of a deep-sea discovery system must therefore be hydrophones, samplers
and cameras in the water, and the satellite's job is to say where to put
them.

The second half is a learning question with an unusual twist. A classifier
trained on known species will confidently assign a new species to the closest
old one; discovery requires the model to say ``none of the above'' --- the
\emph{open-set} recognition problem \cite{hendrycks2017baseline,
lee2018simple, liu2020energy, sun2022knn}. Since no public recording of
undescribed deep-sea animals exists (by definition), we build a simulator of a
deep-sea soundscape: species are parametric call families with
species-specific parameter distributions, propagated to a hydrophone with the
physics of Section~\ref{sec:budget}. We train on $K$ known species, hold out
$U$ unknown ones, and ask six novelty scores to separate them; we then cluster
the clips flagged as novel and ask whether the clusters correspond to the
unknown species --- the step from ``something new'' to ``how many new
things''. To keep the simulator honest we run the same pipeline on real
marine-mammal recordings from the Watkins database in the BEANS packaging
\cite{hagiwara2022beans}.

The last section closes the loop with survey statistics: given detections,
how many species remain unseen (Chao1, Good--Turing \cite{chao1984nonparametric,
good1953population}), and how should a fixed budget of sensor-hours be spread
over sites when a satellite-derived habitat prior is only sometimes right? The
paper's contribution is thus a quantified map --- what can be sensed, what can
be recognised as new, and where to look --- rather than a claim of discovery;
every number regenerates from one command.
